%==============================================================================
% Section IV: Experiments & Results
%==============================================================================
\section{Experiments}
\label{sec:experiments}

We evaluate [METHOD NAME] across multiple benchmarks to answer the
following research questions:
\begin{description}
  \item[\textbf{RQ1}] Does [METHOD NAME] outperform state-of-the-art
                     baselines on standard benchmarks?
  \item[\textbf{RQ2}] How does each component contribute to overall
                     performance?
  \item[\textbf{RQ3}] How sensitive is the method to the choice of
                     hyperparameters?
  \item[\textbf{RQ4}] What is the computational cost compared to existing
                     methods?
\end{description}

% ---- Experimental setup ----------------------------------------------------
\subsection{Experimental Setup}
\label{subsec:setup}

\subsubsection{Datasets}
We evaluate on three public datasets:
\textbf{[Dataset~1]}~\cite{he2016deep} containing \(N_1\) samples,
\textbf{[Dataset~2]}~\cite{devlin2019bert} with \(N_2\) samples, and
\textbf{[Dataset~3]} (proprietary). Standard train/val/test splits are
used.

\subsubsection{Baselines}
We compare against [Baseline~1]~\cite{he2016deep},
[Baseline~2]~\cite{vaswani2017attention}, and
[Baseline~3]~\cite{devlin2019bert}, covering the most competitive
approaches in the literature.

\subsubsection{Evaluation Metrics}
We report \textbf{Accuracy}, \textbf{F1-score}, and \textbf{AUROC} for
classification, and additionally measure \textbf{Inference Latency} (ms)
and \textbf{Peak GPU Memory} (GB) for efficiency.

\subsubsection{Implementation Details}
All models are implemented in PyTorch~1.13 and trained on a single
NVIDIA~A100 (40\,GB) GPU. We use the AdamW optimizer
(\(\eta=1\!\times\!10^{-4}\), weight decay \(0.01\)), batch size~\(128\),
and train for~\(100\) epochs. Hyperparameters are tuned on the validation
set; \(\tau=0.07\) and \(\lambda=0.5\) yielded the best results.

% ---- Main results ----------------------------------------------------------
\subsection{Main Results (RQ1)}
\label{subsec:main-results}

\Cref{tab:main} summarizes our main results. [METHOD NAME] consistently
outperforms all baselines across the three datasets. On~[Dataset~1], we
improve top-1 accuracy by \(+2.4\%\) over the strongest baseline.

\begin{table}[t]
  \caption{Comparison with state-of-the-art on three benchmarks.
           Best results are in \textbf{bold}.}
  \label{tab:main}
  \centering
  \small
  \begin{tabular}{lccc}
    \toprule
    Method                & Dataset~1 & Dataset~2 & Dataset~3 \\
    \midrule
    Baseline~1~\cite{he2016deep}        & 86.1 & 78.3 & 71.2 \\
    Baseline~2~\cite{vaswani2017attention} & 88.7 & 81.0 & 74.5 \\
    Baseline~3~\cite{devlin2019bert}    & 89.5 & 82.4 & 75.9 \\
    \midrule
    \textbf{[METHOD NAME] (ours)} & \textbf{91.9} & \textbf{85.1} & \textbf{78.6} \\
    \bottomrule
  \end{tabular}
\end{table}

% ---- Ablation studies -----------------------------------------------------
\subsection{Ablation Study (RQ2)}
\label{subsec:ablation}

We ablate each component to isolate its contribution
(\Cref{tab:ablation}). Removing [Component~2] degrades accuracy by
\(-1.8\%\), confirming its importance.

\begin{table}[t]
  \caption{Ablation study on Dataset~1.}
  \label{tab:ablation}
  \centering
  \small
  \begin{tabular}{lc}
    \toprule
    Variant                                & Accuracy (\%) \\
    \midrule
    Full model                             & \textbf{91.9} \\
    \quad w/o [Component~1]                & 90.4 \\
    \quad w/o [Component~2]                & 90.1 \\
    \quad w/o [Component~3]                & 89.7 \\
    \bottomrule
  \end{tabular}
\end{table}

% ---- Hyperparameter sensitivity --------------------------------------------
\subsection{Hyperparameter Sensitivity (RQ3)}
\label{subsec:hyperparam}

\Cref{fig:sensitivity} shows that performance is stable across a wide
range of \(\lambda\in[0.1, 1.0]\), indicating robustness to this
hyperparameter.

% \begin{figure}[t]
%   \centering
%   \includegraphics[width=\columnwidth]{sensitivity.pdf}
%   \caption{Accuracy on Dataset~1 vs.\ trade-off coefficient \(\lambda\).}
%   \label{fig:sensitivity}
% \end{figure}

% ---- Efficiency analysis ---------------------------------------------------
\subsection{Efficiency Analysis (RQ4)}
\label{subsec:efficiency}

\Cref{tab:efficiency} compares inference latency and memory footprint.
[METHOD NAME] is~\(1.3\times\) faster than Baseline~3 while using~\(22\%\)
less GPU memory.

\begin{table}[t]
  \caption{Inference latency and peak GPU memory.}
  \label{tab:efficiency}
  \centering
  \small
  \begin{tabular}{lcc}
    \toprule
    Method      & Latency (ms) \(\downarrow\) & Memory (GB) \(\downarrow\) \\
    \midrule
    Baseline~3  & 24.6 & 6.8 \\
    \textbf{Ours} & \textbf{18.9} & \textbf{5.3} \\
    \bottomrule
  \end{tabular}
\end{table}

% ---- Discussion / qualitative results --------------------------------------
\subsection{Discussion}
\label{subsec:discussion}

The gains of [METHOD NAME] are most pronounced on [hard subset], where
prior methods struggle due to [reason]. Qualitative visualizations
(\Cref{fig:qualitative}) further confirm that [observation].
Nonetheless, we observe that the method underperforms when
[failure mode] --- we discuss potential remedies in
\Cref{sec:conclusion}.
